% ==================================================================
\section{Iterative: Fine Grain Task Decomposition}
\label{sec:iter-fine}
% ==================================================================

\subsection{Implementation}
\label{sec:iter-fine-impl}

The \textbf{fine grain iterative} strategy extends the Tile version by exploiting the
additional parallelism that exists \emph{inside} each tile.
Rather than completing the border check and fill/compute sequentially within each tile
task, the fine grain version exposes three independent sub-tasks per tile:

\begin{enumerate}
    \item A task that computes the \emph{vertical} borders of the tile (left and right
          columns), setting \texttt{eV}.
    \item A task that computes the \emph{horizontal} borders (top and bottom rows,
          including all four corners), setting \texttt{eH}.
    \item A task (spawned after \texttt{\#pragma~omp~taskgroup} ensures both border tasks
          have finished) that either fills the tile uniformly or, in the non-uniform case,
          spawns one child task per row to compute all interior pixels.
\end{enumerate}

The outer tile task from the Tile version is \textbf{kept} as a wrapper; the three inner
tasks are created inside it.

\begin{lstlisting}[caption={Key parallel section of \texttt{fine-grain.c} -- inner tasks
                   for each tile.}, label={lst:fine-grain}]
#pragma omp parallel
#pragma omp single
for (int y = 0; y < ROWS; y += TILE)
    for (int x = 0; x < COLS; x += TILE) {
        #pragma omp task firstprivate(x, y)
        {
            int eV, eH;

            #pragma omp taskgroup
            {
                // Task 1: vertical borders (left & right columns)
                #pragma omp task shared(eV)
                {
                    eV = 1;
                    int ref = pixel_dwell(..., x, y, ...);
                    for (int py = y; py < y + TILE; py++) {
                        M[py][x]        = pixel_dwell(..., x,        py, ...);
                        M[py][x+TILE-1] = pixel_dwell(..., x+TILE-1, py, ...);
                        eV &= (ref == M[py][x]);
                        eV &= (ref == M[py][x+TILE-1]);
                    }
                }
                // Task 2: horizontal borders (top & bottom rows)
                #pragma omp task shared(eH)
                {
                    eH = 1;
                    for (int px = x; px < x + TILE; px++) {
                        M[y][px]        = pixel_dwell(..., px, y,        ...);
                        M[y+TILE-1][px] = pixel_dwell(..., px, y+TILE-1, ...);
                        eH &= (M[y][x] == M[y][px]);
                        eH &= (M[y][x] == M[y+TILE-1][px]);
                    }
                }
            } // implicit taskwait -- eH and eV are ready here

            // Task 3: fill or compute (spawns per-row child tasks)
            #pragma omp task
            {
                if (eH & eV) {                       // uniform tile
                    for (int py = y; py < y + TILE; py++)
                        #pragma omp task
                        fill_row(M, py, x, TILE, M[y][x], ...);
                } else {                             // non-uniform tile
                    for (int py = y; py < y + TILE; py++)
                        #pragma omp task
                        compute_row(M, py, x, TILE, ...);
                }
            }
        }
    }
\end{lstlisting}

The \texttt{\#pragma~omp~taskgroup} enforces that the result tasks begin only after
\emph{both} border tasks complete, because the fill/compute decision depends on
\texttt{eV} and \texttt{eH}.
Histogram increments remain protected by \texttt{\#pragma~omp~atomic}.
Display updates use \texttt{\#pragma~omp~critical} to serialize X11 calls.

\paragraph{Comparison with Iterative Tile.}
The Tile version creates $\bigl(\mathit{ROWS}/\mathit{TILE}\bigr)^{2}$ tasks in total.
The Fine Grain version creates up to $3 + 2\cdot\mathit{TILE}$ tasks \emph{per tile},
increasing the task count by roughly $2\cdot\mathit{TILE}$ fold.
This exposes more parallelism at the cost of higher task-creation and synchronisation
overhead; whether the trade-off is beneficial depends on tile size and thread count.

\subsection{Modelfactor Analysis}
\label{sec:iter-fine-mf}

% TODO: Run  sbatch submit-strong-extrae.sh mandel-omp-fine-grain
%       Paste Modelfactors tables and compare with the Tile version.

\todo[inline]{Insert Modelfactors table for \texttt{fine-grain} (16 threads).
Compare with Iterative Tile: does finer granularity improve useful parallelism or does
scheduling overhead dominate?}

\subsection{Paraver Analysis}
\label{sec:iter-fine-paraver}

% TODO: Include Paraver views for 16 threads and compare with Iterative Tile.
%   - "Useful/Instantaneous parallelism"
%   - "OpenMP tasking / explicit tasks function created and executed"
%   Compare task density, idle gaps, and load balance with the Tile trace.

\todo[inline]{Insert Paraver screenshots (16 threads) for Fine Grain and compare with
the Tile trace. Comment on whether the extra tasks reduce idle time or introduce
additional overhead.}

% \begin{figure}[H]
%   \centering
%   \begin{subfigure}[b]{0.48\linewidth}
%     \includegraphics[width=\linewidth]{fotos/iter-tile-paraver-tasks.png}
%     \caption{Iterative Tile.}
%   \end{subfigure}
%   \hfill
%   \begin{subfigure}[b]{0.48\linewidth}
%     \includegraphics[width=\linewidth]{fotos/iter-fine-paraver-tasks.png}
%     \caption{Iterative Fine Grain.}
%   \end{subfigure}
%   \caption{Task creation view -- 16 threads.}
%   \label{fig:fine-paraver-compare}
% \end{figure}

\subsection{Strong Scalability}
\label{sec:iter-fine-scalability}

% TODO: Include speedup plot and compare with Iterative Tile.

\todo[inline]{Insert the strong scalability plot for Fine Grain and compare with
the Tile speedup curve. Discuss whether additional task granularity improves
or hurts scalability at high thread counts.}

% \begin{figure}[H]
%   \centering
%   \includegraphics[width=0.75\linewidth]{fotos/iter-fine-speedup.png}
%   \caption{Strong scalability comparison: Iterative Tile vs.\ Fine Grain.}
%   \label{fig:fine-speedup}
% \end{figure}
